\documentclass{article}
\usepackage{amsmath}
\usepackage{amsthm}
\usepackage{amssymb}
\usepackage{geometry}
\geometry{a4paper, margin=1in}

\theoremstyle{plain}
\newtheorem{lemma}{Lemma}

\usepackage{graphicx}
\usepackage{caption}

\title{One-off Binary Tree}
\author{Ivo Knittel}
\date{11.1.2025}

\begin{document}

\maketitle

\begin{abstract}
This document introduces the concept of a one-off binary tree, focusing on entities termed "Mergables". These entities are characterized by specific properties, qualities, and ranges. A merging operation is defined for these entities, enabling the construction of a binary tree. Applications and theoretical insights into the properties of such a structure are discussed.
\end{abstract}

\tableofcontents

\section{Binary Tree of Mergables}

\subsection*{Mergable}
\subsubsection*{Data Stucture}
An entity of type Mergable has:
\begin{enumerate}
    \item Properties 
    \item A quality in $[0,1]$, where $0$ represents the worst quality and $1$ represents perfect quality
    \item A linear range $[\text{lo}, \text{hi}]$ in $[0, 2^M]$, with $\text{lo} \leq \text{hi}$
\end{enumerate}

 The width of the item $w =\text{hi} -  \text{lo}$.

\subsubsection*{Neighbors}

Two Mergables $v, w$ are neighbors if:
\[
\text{hi}(v) + 1 = \text{lo}(w)
\]

\subsubsection*{Merge}
There exists a merge function:
\[
\text{merge}: v, w  \text{  Mergable}, \text{ and neighbors} \rightarrow u  \text{  Mergable}
\]

with  $\text{lo}(u)=\text{lo}(v)$ and $\text{hi}(u)=\text{hi}(w)$

The merge is associative:
\[
\text{merge} (v_1, \text{merge}(v_2, v_3)) =  \text{merge} (\text{merge}(v_1, v_2), v_3) 
\]

\subsection*{The Binary Tree of Mergables}

With mergables:
\begin{align*}
    v^0_j, & \quad j = 0, \dots 2^M - 1
\end{align*}

there are also:
\begin{align*}
    v^1_j, & \quad j = 0, \dots 2^{M-1} - 1
\end{align*}

with:
\[
v^1_j = \text{merge}(v^0_{2j}, v^0_{2j+1})
\]

and end index $e^g = 2^{M-g} - 1$.

In general, for $g = 1, \dots M$:
\begin{align*}
    v^{g+1}_j, & \quad j = 0, \dots 2^{M-g} - 1
\end{align*}

with:
\[
v^{g+1}_j = \text{merge}(v^g_{2j}, v^g_{2j+1})
\]

\subsection*{Example}

\subsubsection*{Properties}
\begin{itemize}
    \item $\Sigma x$: Real number
    \item $\Sigma x^2$: Real number
    \item $n$: Integer
    \item $\text{quality}$: Real number
\end{itemize}

For $v^0_j$, the properties are:
\begin{align*}
    \Sigma x & = v^0_j \\
    \Sigma x^2 & = (v^0_j)^2 \\
    n & = 1 \\
    \text{quality} & = \text{undefined}
\end{align*}

\subsubsection*{Merge Function}

For $\text{merge}(v, w) \rightarrow u$:
\begin{align*}
    u_{\Sigma x} & = v_{\Sigma x} + w_{\Sigma x} \\
    u_{\Sigma x^2} & = v_{\Sigma x^2} + w_{\Sigma x^2} \\
    \text{error} & = (\Sigma x^2 / n - (\Sigma x / n)^2) / (\Sigma x / n)^2
\end{align*}

\subsubsection*{Best Quality (Maximum Equality)}
 $x_j \quad j = 0...(n-1)$ with $x_j = 1$:
\begin{align*}
    \Sigma x &  = n \\
    \Sigma x^2 & = n \\
    \Sigma x^2 / n & = 1 \\
    \Sigma x / n & = 1 \\
    \text{error} & = 0
\end{align*}

\subsubsection*{Worst Quality (Maximum Inequality)}
 $x_j \quad j = 0...(n-1)$ with $x_0 = 1$ , else $x_j = 0 $:
\begin{align*}
    \Sigma x &  = n \\
    \Sigma x^2 & = n^2 \\
    \Sigma x / n & = 1 \\
    \Sigma x^2 / n & = n \\
    \text{error} & = n - 1 \\
\end{align*}

\subsubsection*{Quality}
     with 
\[ 
\text{quality}  = ((n-1) - \text{error}) / (n-1),
\]
the best quality value is indeed $=1$, and the worst quality $=0$.
\section{One-off Binary Tree}

The properties of the One-off Binary Tree are
\begin{itemize}
    \item The overall qualities in the binary tree can be improved by some freedom in the merging of items.
    \item The binary tree topology is preserved, to leverage the algorithms based on it.
\end{itemize}

\subsection*{One-off shifts}
At the base of that tree are the same Mergables as before:
\begin{align*}
    v^0_j, & \quad j = 0, \dots e^0
\end{align*}

For the One-Off-Binary tree, with a vector of ternary numbers 
\begin{align*}
   t^1_{0} = 0, &   t^1_{j} , t \in {-1, 0, 1},  &  j = 0, \dots e^0
\end{align*}

ranges in the first row $g=1$ are

\begin{align*}
   \text{lo}( v^1_j) = 2j +t^1_j , &  \text{hi}( v^1_j) = 2j+1+t^1_(j+1)  &  \quad j = 0, \dots e^1
\end{align*}

\subsubsection*{Size limits}
On the shift vector of on-off shifts  ${t}$, the condition
\[
\left| t^1_{j+1} - t^1_j \right| \le 1
\]

is imposed, meaning only item sizes of 1, 2, or 3 are allowed.  
Items are not allowed to vanish, because this would change the
tree topology. items of size 4 can be replaced by two pairs, which are then merged
on the next-higher row.

\subsubsection*{One-off tree properties}
For the second row $g=2$, with its own vector of ternaries 

\begin{align*}
   t^2_{0} = 0, &   t^2_{j} , t \in {-1, 0, 1},  &  j = 0, \dots e^2
\end{align*}

ranges are

\begin{align*}
   \text{lo}( v^2_j) = \text{lo}(v^1_{2 j + t^2_j})  , &  \text{hi}( v^2_j) =\text{hi}(v^1_{2(j+1)+1+t^2_{j+1}})   &  \quad j = 0, \dots, e^2
\end{align*}

For any row $g>0$, ranges are

\begin{align*}
   \text{lo}( v^g_j) = \text{lo}(v^{g-1}_{2 j + t^g_j})  , &  \text{hi}( v^g_j) =\text{hi}(v^{g-1}_{2(j+1)+1+t^g_{j+1}})   &  \quad j = 0, \dots, e^g
\end{align*}

\begin{lemma}
For the width of a Mergable in row $g$ of a One-Off-Binary tree all values are possible from $1$  to $ e^g$.
\end{lemma}

\begin{proof}
(left to the reader)
\end{proof}

\subsection*{ Merge Candidates}

In the binary tree, for any item, the merge is determined:
\[
\begin{array}{c|ccccc}
\text{item} & v_0 & v_1 & v_2 & v_3 & v_4 \\
\hline
\text{merge} & m(v_0, v_1) & m(v_1, v_2) & m(v_2, v_3) & m(v_3, v_4) & m(v_4, v_5)
\end{array}
\]

In the one-off binary tree, for any item candidate merges are generated to be processed further:
\[
\begin{array}{c|ccccc}
\text{item} & v_0 & v_1 & v_2 & v_3 & v_4 \\
\hline
\text{candidate\_lo} &   & m(v_0, v_1) & m(v_1, v_2) & m(v_2, v_3) & m(v_3, v_4)  \\
\text{candidate\_hi} & m(v_0, v_1) & m(v_1, v_2) & m(v_2, v_3) & m(v_3, v_4) & m(v_4, v_5)
\end{array}
\]

\subsection*{ Merge Candidate Preference}

Item $v_1$ has two candidate merges and chooses the better one (marked by "!" in the table), the same with $v_2$.
If their choices agree, and we get a "happy couple": 
\[
\begin{array}{c|cc}
\text{item}  & v_1 & v_2 \\
\hline
\text{candidate\_lo} & m(v_0, v_1) & m(v_1, v_2)! \\
\text{candidate\_hi} & m(v_1, v_2)! & m(v_2, v_3) 
\end{array}
\]

In the other case we get a "a midsummer night's dream".
$v_1$ loves $v_2$, but $v_2$ loves $v_3$ etc. 
\[
\begin{array}{c|cc}
\text{item}  & v_1 & v_2 \\
\hline
\text{candidate\_lo} & m(v_0, v_1) & m(v_1, v_2) \\
\text{candidate\_hi} & m(v_1, v_2)! & m(v_2, v_3)! 
\end{array}
\]



\subsubsection{Triples}
Assuming for now qualities cannot be equal,
After settling what is left beside couples are
chains of one element, Triples. 
A Triple is a single, attracted to one pair. 
If it is the lower-index side, denoted as $T^{-}$:

\[
\begin{array}{c|ccccccc}
\text{item}  & 1 & 2 & 3 & 4 & 5\\
\hline
\text{candidate\_lo} & (0, 1) & (1, 2)! & (2, 3)! & (3, 4)!  & (4, 5)! \\
\text{candidate\_hi} & (1, 2)! & (2, 3) & (3, 4) & (4, 5)  & (5, 6)
\end{array}
\]

If the pair is the upper-index side, denoted as  $T^{+}$:

\[
\begin{array}{c|ccccccc}
\text{item}  & 1 & 2 & 3 & 4 & 5\\
\hline
\text{candidate\_lo} & (0, 1) & (1, 2)! & (2, 3) & (3, 4)  & (4, 5)! \\
\text{candidate\_hi} & (1, 2)! & (2, 3) & (3, 4)! & (4, 5)!  & (5, 6)
\end{array}
\]

After settling, a tree row $g$ consists of triples $T^\pm_{ k} , k =0...N_T, N_T < e^g/3$.
Between triples can exist any number of Couples, or none.

\subsection*{Preserving tree topology}
We now decide what to do with pairs of Triples,  $T^{+}  T^{+}, T^{+}  T^{-} , T^{-}  T^{+}$, and $T^{-}  T^{-}$, without Couples in-between.
Figures 1-2 recall the basic property of the One-Off Binary Tree. The boundaries between
the $j^0$ items to be merged into $j^1$ items are allowed to move by one unit.

If all $j^0$ are Couples, the case  $t^1_j = 0$, is the original Binary Tree.
 The $j^0_j$ items are merged as in Fig. 1, or shortanhd as 
\[
| 0 1 |2  3| 4 5 | ... | (e^0-1)e^0 |.   
\]
For $t^1_j = 1$,  items are merged according to 
\[
 | 0 1 2 |  3 4 | 5 6 | ... | (e^0-1) | e^0 |  
\]
For $t^1_j = -1$,  items are merged according to
\[
  | 0  |  1 2 |  3 4 | ...   | (e^0-2)(e^0-1)e^0 | | 
\]

\begin{figure}[h]
    \centering
    \resizebox{\textwidth}{!}{\includegraphics{MergableGraph1.pdf}}
    \caption{The base row $j_0$ and the first row $j_1$ of a One-Off-Binary-Tree with all unit shifts $t_j=0$. In each colunm, the $j_1$ item is the reult of merge of the  $j_0$ items.}
    \label{fig:tikz_diagram}
\end{figure}

\begin{figure}[h]
    \centering
    \resizebox{\textwidth}{!}{\includegraphics{MergableGraph2.pdf}}
    \caption{Setting $t^1_0 = 1$ means the node 2 being merged into node 0 instead of node 1.}
    \label{fig:tikz_diagram}
\end{figure}


\subsubsection{Triple pairs}


\[
\begin{array}{c|cccccccccccc}
  & 1 & 2 & 3 & 4 & 5& 6 & 7 & 8 & 9 & 10\\
\hline
& (0, 1) & (1, 2)! & (2, 3) & (3, 4)  & (4, 5)! & (5, 6) & (7, 8)! & (8, 9) & (9, 10)  & (10, 11)! \\
 & (1, 2)! & (2, 3) & (3, 4)! & (4, 5)!  & (5, 6) & (6, 7)! & (8, 9) & (9, 10)! & (10, 11)!  & (11, 12)
\end{array}
\]

\section*{Conclusions}

The one-off binary tree demonstrates the utility of merging entities with specific properties and qualities. By defining a merge operation and constructing a binary tree, the structure supports analysis of data aggregation and quality metrics. Future work could explore extensions to other data types and applications.


end{document}